\chapter{Đúng lý, sai tình}
\markboth{Đúng lý, sai tình}{Right in Logic, Wrong in Feeling}

\section{Một người nói thật nhưng nói đúng lúc người khác đang đau.}
\begin{truyen}{Một người nói thật nhưng nói đúng lúc người khác đang đau.}{Truth without timing can wound.}
\chuhoa{A}{nh không sai về nội dung, nhưng sai ở thời điểm và cách nói. Khi người đối diện còn đang tổn thương, sự thật khô cứng nghe như thêm một nhát cắt. Đúng về lý mà thiếu tình vẫn có thể làm sai điều cần giữ.}
\textit{(He was not wrong in content, but he was wrong in timing and tone. When the other person was still hurting, hard truth sounded like another cut. Logical correctness without tenderness can still damage what matters.)}
\end{truyen}
\begin{baihoc}
Sự thật cũng cần lòng trắc ẩn để đến đúng nơi trong lòng người khác.
\textit{(Truth also needs compassion to arrive well in another person's heart.)}
\end{baihoc}
\begin{ghinhoanh}\textbf{Short English to remember:}

Truth without timing can wound.
\end{ghinhoanh}
\begin{tuhocnhanh}
1. Em có đang đúng mà nói quá mạnh không? \textit{(Are you right but speaking too harshly?)}

2. Điều gì cần được nói lại bằng một cách mềm hơn? \textit{(What needs to be said again more gently?)}
\end{tuhocnhanh}
\ngancach

\section{Một người trách cha mẹ bằng lý lẽ đúng nhưng giọng đầy khinh.}
\begin{truyen}{Một người trách cha mẹ bằng lý lẽ đúng nhưng giọng đầy khinh.}{Correct words can carry the wrong spirit.}
\chuhoa{C}{ậu có những điều nói ra không sai: cha mẹ có thiếu sót, có lúc áp đặt. Nhưng cách cậu nói làm mọi điều đúng trở nên cay đắng. Cậu thắng ở luận điểm, nhưng làm gia đình thua ở sự gần gũi.}
\textit{(He was not wrong: his parents had flaws and could be controlling. But the way he spoke made every right point taste bitter. He won the argument, yet the family lost closeness.)}
\end{truyen}
\begin{baihoc}
Điều đúng mà đi cùng thái độ coi thường thường không chữa được gì.
\textit{(What is right, when carried with contempt, rarely heals anything.)}
\end{baihoc}
\begin{ghinhoanh}\textbf{Short English to remember:}

Correct words can carry the wrong spirit.
\end{ghinhoanh}
\begin{tuhocnhanh}
1. Em có đang dùng điều đúng như vũ khí không? \textit{(Are you using what is right as a weapon?)}

2. Em có thể nói điều đó bằng thái độ khác nào? \textit{(How could you say it with a different spirit?)}
\end{tuhocnhanh}
\ngancach

\section{Một giáo viên phê bình đúng lỗi nhưng làm trò mất mặt.}
\begin{truyen}{Một giáo viên phê bình đúng lỗi nhưng làm trò mất mặt.}{Correction also needs dignity.}
\chuhoa{T}{rò làm sai thật và cần được nhắc. Nhưng khi lời nhắc được nói theo cách làm em xấu hổ trước đám đông, bài học đúng bị phủ lên bằng nỗi tủi thân.}
\textit{(The student was truly wrong and needed correction. But when the correction was delivered in a humiliating way before others, the right lesson was covered by shame.)}
\end{truyen}
\begin{baihoc}
Sửa người khác cho đúng mà không giữ phẩm giá cho họ thì cái đúng ấy còn thiếu một nửa.
\textit{(Correcting someone rightly without preserving their dignity leaves that rightness incomplete.)}
\end{baihoc}
\begin{ghinhoanh}\textbf{Short English to remember:}

Correction also needs dignity.
\end{ghinhoanh}
\begin{tuhocnhanh}
1. Khi góp ý người khác, em có giữ thể diện cho họ không? \textit{(When you correct others, do you protect their dignity?)}

2. Em từng tổn thương vì bị góp ý kiểu nào? \textit{(What kind of correction has hurt you before?)}
\end{tuhocnhanh}
\ngancach

\section{Một người giúp tiền đúng lúc nhưng nhắc ơn mãi về sau.}
\begin{truyen}{Một người giúp tiền đúng lúc nhưng nhắc ơn mãi về sau.}{Help can turn wrong when it humiliates.}
\chuhoa{B}{an đầu, sự giúp đỡ là thật. Nhưng khi nó đi kèm những lần nhắc đi nhắc lại và thái độ bề trên, người nhận không còn thấy được cứu mà thấy bị giữ nợ tinh thần mãi mãi.}
\textit{(At first, the help was real. But when it came with repeated reminders and superiority, the receiver no longer felt rescued, only trapped in emotional debt.)}
\end{truyen}
\begin{baihoc}
Điều tốt nếu làm để giữ quyền hơn người thì sẽ dần méo đi.
\textit{(A good deed done to maintain superiority gradually becomes distorted.)}
\end{baihoc}
\begin{ghinhoanh}\textbf{Short English to remember:}

Help turns wrong when it humiliates.
\end{ghinhoanh}
\begin{tuhocnhanh}
1. Khi giúp ai, em có mong quyền lực gì trên họ không? \textit{(When helping someone, do you expect any power over them?)}

2. Em muốn được giúp theo cách nào nếu mình rơi vào khó khăn? \textit{(How would you want to be helped if you were in difficulty?)}
\end{tuhocnhanh}
\ngancach

\section{Một người quyết định đúng cho con nhưng không hề lắng nghe con.}
\begin{truyen}{Một người quyết định đúng cho con nhưng không hề lắng nghe con.}{Right choices need listening too.}
\chuhoa{C}{ha mẹ có thể nhìn xa hơn và muốn điều tốt hơn cho con. Nhưng nếu mọi điều tốt đều áp xuống mà không có lắng nghe, đứa trẻ chỉ học được sự phục tùng hoặc phản kháng, chứ chưa chắc học được điều đúng.}
\textit{(Parents may see further and want what is best for a child. But if all good decisions are imposed without listening, the child learns obedience or resistance, not necessarily wisdom.)}
\end{truyen}
\begin{baihoc}
Đúng cho người khác mà thiếu lắng nghe người đó thì cái đúng rất dễ thành áp đặt.
\textit{(Being right for someone else without listening to them easily turns rightness into control.)}
\end{baihoc}
\begin{ghinhoanh}\textbf{Short English to remember:}

Right choices need listening too.
\end{ghinhoanh}
\begin{tuhocnhanh}
1. Em có đang muốn tốt cho ai theo cách quá áp đặt không? \textit{(Are you wanting good for someone in a way that is too controlling?)}

2. Em cần nghe thêm điều gì từ họ? \textit{(What do you need to hear more from them?)}
\end{tuhocnhanh}
\ngancach